% ==============================================================================
%  Comandos / notación común a todos los temas
%  (usar SIEMPRE estos comandos para que la notación sea homogénea)
% ==============================================================================

% --- Conjuntos y probabilidad ---------------------------------------------------
\newcommand{\R}{\mathbb{R}}
\DeclareMathOperator{\E}{E}                  % esperanza      \E[Y]
\DeclareMathOperator{\Var}{Var}              % varianza       \Var(Y)
\DeclareMathOperator{\Cov}{Cov}              % covarianza     \Cov(X,Y)
\DeclareMathOperator{\Cor}{Cor}              % correlación
\newcommand{\Normal}{\mathcal{N}}            % \Normal(0,\sigma^2)
\newcommand{\iid}{\overset{\text{i.i.d.}}{\sim}}
\newcommand{\bajoHo}{\overset{H_0}{\sim}}    % distribución bajo H0

% --- Vectores y matrices (negrita) ----------------------------------------------
\newcommand{\vect}[1]{\bm{#1}}               % \vect{Y}, \vect{\beta}
\newcommand{\T}{^{\mathsf{T}}}               % traspuesta: X\T

% --- Sumas de cuadrados y medias cuadráticas ------------------------------------
\newcommand{\SC}[1]{SS_{\text{#1}}}          % \SC{hip}   -> SS_hip
\newcommand{\MC}[1]{MS_{\text{#1}}}          % \MC{res}   -> MS_res

% --- Estadística descriptiva --------------------------------------------------
\newcommand{\media}[1]{\overline{#1}}        % \media{y}
\newcommand{\est}[1]{\widehat{#1}}           % estimador: \est{\mu}

% --- Distribuciones -----------------------------------------------------------
\newcommand{\Fdist}[2]{F_{#1,\,#2}}          % \Fdist{1}{n-1}
\newcommand{\tdist}[1]{t_{#1}}
\newcommand{\chisq}[1]{\chi^2_{#1}}

% --- Texto --------------------------------------------------------------------
\newcommand{\Rsoft}{\textsf{R}}              % el software R
\newcommand{\codigo}[1]{\texttt{#1}}

% --- Otros ----------------------------------------------------------------------
\newcommand{\cond}{\mid}                     % condicionado: \E(Y \cond X=x)
\newcommand{\gl}{\text{g.l.}}                % grados de libertad
\newcommand{\sumi}{\sum_{i=1}^{n}}           % suma habitual de 1 a n

% --- Regresión (Tema 3 en adelante) ----------------------------------------------
\newcommand{\Ynew}{Y_{h(\mathrm{new})}}      % nueva observación en x_h
\newcommand{\Xnew}{X_{h(\mathrm{new})}}      % predicción inversa
\newcommand{\Ymnew}{\media{Y}_{h(\mathrm{new})}} % media de m nuevas observaciones
\newcommand{\tq}[2]{t_{#1;\,#2}}             % cuantil: \tq{n-2}{1-\alpha/2}
\newcommand{\Fq}[3]{F_{#1,\,#2;\,#3}}        % cuantil: \Fq{1}{n-2}{1-\alpha}
\newcommand{\ee}{\operatorname{ee}}          % error estándar
\newcommand{\SSX}{SS_X}
\newcommand{\raizMSE}[1]{\sqrt{MSE\left(#1\right)}}

% --- Validación (Tema 4 en adelante) -------------------------------------------
\newcommand{\DFFits}{\mathrm{DFFits}}        % \DFFits_i
\newcommand{\DFBeta}{\mathrm{DFBeta}}        % \DFBeta_{ji}
\newcommand{\sumk}{\sum_{k=1}^{n}}           % suma en k de 1 a n

% --- Notación matricial (Tema 5 en adelante) ------------------------------------
\newcommand{\XtX}{\vect{X}\T\vect{X}}                 % X'X
\newcommand{\XtXi}{(\vect{X}\T\vect{X})^{-1}}         % (X'X)^{-1}
\newcommand{\bhat}{\est{\vect{\beta}}}                % beta estimado (vector)
